\chapter{Electromagnetic Induction}

\section{Faraday's Law and Lenz's Law}

\textbf{Faraday's law} and \textbf{Lenz's law} are usually written together as
\begin{equation} \label{eq:faraday-lenz-law}
    \epsilon = - \odv{\Phi_B}{t}
\end{equation}
Strictly speaking, Faraday's law gives the magnitude,
while Lenz's law gives the sign.

This should not be surprising at this point,
but the direction of the induced emf is also based on the right-hand rule.
Point your right thumb in the direction that the flux is \emph{changing}.
If the magnitude of the flux is \emph{increasing},
point your right thumb in the \emph{same} direction as the magnetic field.
If the magnitude of the flux is \emph{decreasing},
point your right thumb in the \emph{opposite} direction of the magnetic field,
or alternatively, point your \emph{left} thumb
and point in the \emph{same} direction as the magnetic field.
In any case,
the induced emf will be in the same direction as your curled fingers.

\section{Inductors}

Physically speaking, \textbf{inductors} are not much different from solenoids;
they are both coils of wire.
The only difference is how they are used.

Circuits are essentially large loops of conducting metal,
which make them also effective at catching changing magnetic flux and
generating an emf.
The generated emf may overload the electrical devices in the circuit
or counteract the power source;
presumably, these are not desired outcomes.
To solve this issue, we add an inductor.
When the current changes,
the inductor will generate a counter-emf to help stabilize the current.

The effectiveness of an inductor is its \textbf{inductance}~(symbol: \(L\)),
which has SI units of henries~(\unit{\henry}) and is defined as
\begin{equation}
    L \coloneq \frac{V}{\odv{I}{t}}.
\end{equation}

\section{LC Circuits}

\textbf{LC circuits}, as the name suggests,
are composed of an inductor and a capacitor.
If we hook up a charged capacitor to an inductor
it forms an oscillation between storing energy in electric and magnetic fields.

The angular frequency~\(\omega\) of the LC circuit is given by
\begin{equation}
    \omega = \frac{1}{\sqrt{L C}}.
\end{equation}

\section{Maxwell's Equations}

\textbf{Maxwell's equations} are essentially the essence
of electricity and magnetism distilled into four equations.
You have more or less seen them already.
The equations are
\footnote{
    The following equations are actually taken from my notes for
    Waves and Oscillations.
}
\begin{subequations}
    \begin{equation}
        \oiint_S \vecb{E} \cdot \odif{\vecb{A}} = \frac{q}{\epsilon_0},
    \end{equation}
    \begin{equation}
        \oiint_S \vecb{B} \cdot \odif{\vecb{A}} = 0,
    \end{equation}
    \begin{equation}
        \oint_C \vecb{E} \cdot \odif{\vecb{\ell}} = - \odv{\Phi_B}{t},
    \end{equation}
    \begin{equation}
        \oint_C \vecb{B} \cdot \odif{\vecb{\ell}}
        = \mu_0 I + \mu_0 \epsilon \odv{\Phi_E}{t}.
    \end{equation}
\end{subequations}
Indeed, these are Gauss's law for electricity~[\pgRef{eq:gausss-law}],
Gauss' law for magnetism~[\pgRef{eq:gausss-law-magnetism}],
the Faraday-Lenz law~[\pgRef{eq:faraday-lenz-law}],
and Ampère's law~[\pgRef{eq:amperes-law}] (sort of).
It turns out that Ampère's law was incomplete,
and Maxwell added the
\begin{equation*}
    \mu_0 \epsilon \odv{\Phi_E}{t}
\end{equation*}
term, which is often uncreatively called "Maxwell's addition to Ampère's law".
The new combined equation is known as the \textbf{Ampère-Maxwell law}.

The Faraday-Lenz law suggests that changing magnetic fields
create electric fields,
and the Ampère-Maxwell law suggests that changing electricity fields
create magnetic fields.
This suggests that it is possible for these two effects to chain together
into oscillating electric and magnetic fields,
known as \textbf{electromagnetic waves} or \textbf{EM waves}.
Maxwell correctly postulated that his mathematical electromagnetic waves
were in fact light.
From Maxwell's equations,
we can derive that the speed of light \(c\) is
\begin{equation}
    c = \frac{1}{\sqrt{\epsilon_0 \mu_0}}
\end{equation}
The exact derivation is too advanced for this volume
and will be provided in Vol.~3: Waves and Oscillations.

The discovery that light has a fixed speed would be crucial to the
discovery of the theory of relativity in the early 1900s.
